\section{Conclusion}

We began by observing that GLUE was replaced by SuperGLUE within twelve months---and that no automated system detected the saturation that made this necessary.
Our work shows that it could have: by formalizing benchmark score-over-time trajectories as logistic growth processes and applying AIC model selection, we achieve decisive statistical justification for the logistic model ($\Delta\text{AIC} \approx -194$ to $-250$ versus linear, $-113$ to $-357$ versus power-law) and retrospective saturation detection within 3--5 months of community-recognized events.
An empirically optimized criterion achieves one-month accuracy.

Our main contributions are: (1) to our knowledge, the first automated benchmark saturation pipeline requiring no new test data; (2) formal information-theoretic validation of the logistic model for benchmark score dynamics; (3) a dual-criterion saturation detector with documented precision and baseline comparison; (4) discovery that logistic inflection points for GLUE and SuperGLUE predate their launch; and (5) boundary condition evidence extending the pipeline to $\geq$30 entries.

Future work should prioritize: HuggingFace API integration for live monitoring; prospective forecasting with shorter lookback windows; validation on real small benchmarks; and extension to vision and multilingual benchmarks.

Returning to our opening: if this pipeline had been running in 2019, it would have flagged GLUE's saturation within three months of the community-recognized event---and the negative $t_0$ finding suggests the signal was present even before the leaderboard opened.
The benchmark was saturated before we started measuring it.
That is, perhaps, the most important thing our pipeline can tell us.
